\section{Related Work}
\label{sec:related}
Clinical de-identification research shows that removal and pseudonymization
can affect downstream tasks differently
\cite{berg2020impact,bencheikhlarbi2023clinical,singh2025redactor}.
Prompt-obfuscation and task-aware anonymization work also measures privacy and
utility together
\cite{pasch2025balancing,lin2025emojiprompt,loiseau2025taueval,loiseau2026adaptive}.
DiagPriv applies this concern to linked operational evidence.

HMAC is a keyed message-authentication construction
\cite{krawczyk1997hmac}. Stable HMAC pseudonyms remove literals but keep
within-key linkage. DiagPriv's format-preserving condition keeps character
classes; it is not the cryptographic FPE defined by NIST
\cite{dworkin2016fpe}.

LLM-assisted operational RCA is established
\cite{ahmed2023recommending,chen2024automatic,wang2024rcagent,xu2025openrca}.
OpenRCA and Cloud-OpsBench provide broader operational benchmarks
\cite{xu2025openrca,wang2026cloudopsbench}. DiagPriv is smaller and synthetic.
Its focus is the evidence representation sent to the same external model.

Temporal benchmarks and contrast sets motivate controlled near-neighbor cases
\cite{chu2024timebench,gardner2020contrast}. Citation research also separates
answer quality from evidence quality
\cite{gao2023citations,zhang2024finegrained}. DiagPriv combines these ideas in
a scoped study. It does not claim to originate anonymization, RCA
benchmarking, temporal reasoning, or evidence-support scoring.
